
%\documentclass[./main.tex(path of main file)]{subfiles}
\documentclass[10pt]{article}
\usepackage{amsmath}
\DeclareMathOperator*{\argmin}{arg\,min} % thin space, limits underneath in displays
\DeclareMathOperator*{\argmax}{arg\,max} % thin space, limits underneath in displays
\newtheorem{thm}{Theorem}
\usepackage{amssymb}
\usepackage{amsfonts}
\usepackage{mathrsfs}
\usepackage{bm}
\usepackage{indentfirst}
\setlength{\parindent}{0em}
\usepackage[margin=1in]{geometry}
\usepackage{graphicx}
\usepackage{setspace}
\doublespacing
\usepackage[flushleft]{threeparttable}
\usepackage{booktabs,caption}
\usepackage{float}
\usepackage[sort,comma]{natbib}
\usepackage[hidelinks]{hyperref}
\usepackage{booktabs}
\usepackage{multirow}

\usepackage{import}
\usepackage{xifthen}
\usepackage{pdfpages}
\usepackage{transparent}
\usepackage{subfiles}
\usepackage{blindtext}
\usepackage{glossaries}
\usepackage{pdflscape}
\usepackage{adjustbox}
\usepackage{rotating}
\usepackage{eurosym,geometry,ulem,subcaption,caption,color,sectsty,comment,footmisc,pdflscape,array}

% Format definition list
\newenvironment{mylist}[1]%
  {\begin{list}{}{\settowidth{\labelwidth}{\textbf{#1}}
   \setlength{\leftmargin}{\labelwidth}
   \addtolength{\leftmargin}{\labelsep}
   \renewcommand{\makelabel}[1]{\textbf{\hfill##1}}}}%
  {\end{list}}


%%% Format the way to reference subplots
% 1. Change the subfigure label to include parentheses: (a), (b), etc.
\renewcommand\thesubfigure{(\alph{subfigure})}

% 2. Fix the label under the image so it doesn't show double parentheses like ((a))
\captionsetup[subfigure]{labelformat=simple}

% 3. Add a space between the Figure number and the subfigure letter in references
\makeatletter
\renewcommand\p@subfigure{\thefigure~}
\makeatother





\newcommand{\incfig}[1]{%
\def\svgwidth{\columnwidth}
\import{./figures/}{#1.pdf_tex}
}




\title{}
\author{}
\date{}


\begin{document}
\maketitle




\section{Two-way clustering}

Cluster the standard error at time level vs. two-way clustering at time and cohort level.

Cohort is defined by individuals birth year, i.e., all individuals who are born in the same year
belongs to the same cohort.


{\textbf {Cluster at time level}}:
\begin{itemize}
\item Allow for correlation of unkown form in residusals across time periods.
\end{itemize}


{\textbf {Cluster at time and cohort level (two-way clustering)}}:
\begin{itemize}
\item Allow for correlation of unkown form in residusals within the same cohort across time periods
		(the 1980 cohort in 2015 and 2020) and for correlation across cohorts within the same time 
		period (everyone in 2020).
\item It also allows correlation of residuals within the exact same time-cohort cell (the 1980 cohort in 2020)
\end{itemize}

{\textbf {Cluster at time-cohort level}}:
\begin{itemize}
\item It ONLY allows correlation of residuals within the exact same time-cohort cell (the 1980 cohort in 2020)
\item It assumes the correlation across cohorts within the same time period to be ZERO (cohort 1980 and cohort
		1990 in 2020 share no common shocks. This is very restrictive!)
\item It assumes the correlation within the same cohort across time periods to be ZERO (cohort 1980 in 2020 is
		completely independent of the 1980 cohort in 2021. This is very restrictive!)
\item 
		{\textbf {In short: Two-way clustering is much more general and conservative, whereas clustering at the time-cohort cell level is much more restrictive.}}
\end{itemize}



{\textbf {Number of Clusters}}:
\begin{itemize}
\item If you have 60 time periods and 50 cohorts,
		\begin{itemize}
		\item Two-way clustering relies on $ G_{time} = 60 $ and $ G_{cohort} = 50 $
		\item Time-cohort clustering creates $ 60  \times 50 = 3000 $ clusters.
		\end{itemize}
\end{itemize}


{\textbf {Magnitudes of SE}}:
\begin{itemize}
\item 
Clustering at the time-cohort cell level will almost always yield much smaller standard errors than two-way clustering. This is because the time-cohort level creates thousands of tiny clusters and falsely assumes that errors across different years or across different cohorts are completely independent.
\end{itemize}


{\textbf {Interpretation of SE under alternative clutering methods}}:

Two-way clustering computes the variance-covariance matrix as:
\begin{equation*}
 \widehat{Var}_{two-way} =  \widehat{Var}_{time} +  \widehat{Var}_{cohort} -  \widehat{Var}_{time \times cohort}.
 \label{eqn:two_way_clustering_math}
\end{equation*}


Consider results in Table \ref{table:alternative_clustering}, The SE is 0.0054 when it is clustered at time level.
It is 50\% larger than the SE clustered at the cohort level (0.0036). The two-way clustered SE is 0.006.
It suggests that
\begin{itemize}
\item 
		Time-series shocks are the dominant source of error dependence
		\begin{itemize}
    \item 
				The primary unobserved correlation in your data comes from aggregate macroeconomic shocks at time level
				(e.g., business cycles, sudden CPI releases, monetary policy announcements) that hit all households at the same time,
				regardless of their age.
    \item 
				Clustering only at the cohort level would severely understate your standard errors, leading to inflated t-statistics and
    		false positives.
		\end{itemize}

\item 
		Cohort persistence provides an incremental, non-trivial variance contribution
		\begin{itemize}
				\item 
						Moving from $ SE_{time} = 0.0054 $ to $ SE_{two-way} = 0.006 $ increases the SE by roughly 11\% (and the variance 
						by about 23\%).
				\item 
						Economic interpretation: While aggregate time shocks dominate, individuals born in the same year still share 
						persistent unobserved characteristics (e.g., formative lifetime economic experiences, life-cycle stage). 
						Two-way clustering successfully absorbs this extra layer of dependence.
		\end{itemize}	
\end{itemize}

{\textbf {Variance Decomposition}}
In Eq. \eqref{eqn:two_way_clustering_math},
\begin{align*}
\widehat{Var}_{time}    &= 0.0054^{2} = 0.00002916\\
\widehat{Var}_{cohort}  &= 0.0036^{2} = 0.00001296\\
\widehat{Var}_{two-way} &= 0.0060^{2} = 0.00003600
\end{align*}
Hence, 
\begin{equation*}
\widehat{Var}_{time  \times cohort} =  \widehat{Var}_{time} +  \widehat{Var}_{cohort} -  \widehat{Var}_{two-way} = 0.00000612
\end{equation*}

\begin{equation*}
SE_{time  \times cohort} = \sqrt { \widehat{Var}_{time  \times cohort}} \approx 0.00247
\end{equation*}
This reveals the hierachy of SEs:
\begin{equation*}
SE_{time  \times cohort} < SE_{cohort} < SE_{time} < SE_{two-way}
\end{equation*}
This hierarchy confirms that your error structure is well-behaved: standard errors grow as you allow for broader and more general structures of cross-sectional and temporal correlation.


\begin{table}[h!]
\begin{center}
\begin{threeparttable}
\caption{Alternative Clustering}
\label{table:alternative_clustering}

\begin{tabular}{lccc}
\hline\hline
 & (1)& (2)& (3)\\
 \hline \\[-1.8ex]


    $ \tau^{1yr} $&  0.0468***& 0.0468***& 0.0468***\\
                  &    (0.006)&  (0.0036)&  (0.0054)\\
                  &           &          &          \\
       No. of Obs.&    188,073&   188,073&   188,073\\
  Pseudo-R-squared&     0.2057&    0.2057&    0.2057\\



\multicolumn{4}{c}{}\\

Clustering& Time-Cohort& Cohort& Time\\


\hline\hline
\end{tabular}

\end{threeparttable}
\end{center}
\end{table}




{\textbf {Clustering at time and cohort vs. at time and age}}
The two different clustering methods allows for a fundamentally different correlation structure across time.

\begin{align*}
 \widehat{Var}_{time,cohort} &=  \widehat{Var}_{time} +  \widehat{Var}_{cohort} -  \widehat{Var}_{time  \times cohort}\\
 \widehat{Var}_{time,age} &=  \widehat{Var}_{time} +  \widehat{Var}_{age} -  \widehat{Var}_{time  \times age}
\end{align*}

\begin{itemize}
\item 
		$ \widehat{Var}_{time  \times cohort} = \widehat{Var}_{time  \times age} $ because age is determined given time and cohort.
\item 
		The difference is $ \widehat{Var}_{cohort} \ne \widehat{Var}_{age} $.
\item 
		Clustering by age
		\begin{itemize}
		\item 
				Groups different generations across time that happen to be at the same life-cycle stage.
		\item 
				Example: It allows the errors of 30-year-olds in 2010 (born in 1980) to be correlated with 30-year-olds in 2020 (born in 1990).
		\item 
				Economic logic: Captures life-cycle stage-specific shocks (e.g., life milestones, early-career transitions, having young children).
		\end{itemize}
\item 
		Which one is preferred?
		\begin{itemize}
		\item Time + Cohort clustering is strongly preferred.
		\item The leading theoretical concern in expectations data is that individuals carry unobserved, persistent biases formed by their lifetime historical experiences. Because those experiences stick with a birth cohort over time, error persistence follows the birth cohort, not age brackets.
		\item Age-specific patterns (life-cycle consumption) are typically controlled for in the regression index via age fixed effects or polynomials, leaving the unobserved persistence in the cohort dimension.
		\end{itemize}
\end{itemize}









\bibliographystyle{plainnat}
\bibliography{my_bib}

\end{document}

